\documentclass[conference]{IEEEtran}
\IEEEoverridecommandlockouts

\usepackage{cite}
\usepackage{amsmath,amssymb,amsfonts}
\usepackage{algorithmic}
\usepackage{graphicx}
\usepackage{textcomp}
\usepackage{xcolor}
\usepackage{booktabs}
\usepackage{multirow}
\usepackage{hyperref}
\usepackage[utf8]{inputenc}
\usepackage[T5]{fontenc}
\usepackage[vietnamese]{babel}
\usepackage{tikz}
\usetikzlibrary{shapes.geometric, arrows.meta, positioning, fit, backgrounds}

\def\BibTeX{{\rm B\kern-.05em{\sc i\kern-.025em b}\kern-.08em
    T\kern-.1667em\lower.7ex\hbox{E}\kern-.125emX}}

\begin{document}

\title{Nhận Dạng Hành Động Người Bằng Khung Xương}

\author{
  \IEEEauthorblockN{Đỗ Nhật Minh}
  \IEEEauthorblockA{Khoa Công nghệ Thông tin\\
  Trường Đại học Khoa học Tự nhiên, ĐHQG-HCM}
  \and
  \IEEEauthorblockN{Lê Anh Duy}
  \IEEEauthorblockA{Khoa Công nghệ Thông tin\\
  Trường Đại học Khoa học Tự nhiên, ĐHQG-HCM}
}

\maketitle

% Đánh số bảng bằng chữ số Ả Rập thay vì La Mã
\renewcommand{\thetable}{\arabic{table}}

\begin{abstract}
Bài báo này trình bày một pipeline đầu cuối (end-to-end) cho bài toán nhận dạng
hành động người, kết hợp YOLOv8-pose để trích xuất khung xương theo thời gian
thực với Mạng Tích Chập Đồ Thị Không-Thời Gian (ST-GCN) để phân loại hành động.
Hệ thống được đánh giá trên tập con 8 lớp bài tập thể dục của Penn Action Dataset
(1.163 video). Đầu tiên, chúng tôi huấn luyện ST-GCN trên tọa độ khớp ground
truth để thiết lập ngưỡng trên, đạt độ chính xác validation \textbf{98,28\%} và
macro F1-score \textbf{0,9814}. Tiếp theo, chúng tôi thay thế keypoint ground truth
bằng keypoint được ước lượng tự động từ YOLOv8n-pose, ánh xạ 17 khớp COCO sang
13 khớp Penn Action và bổ sung một node trung tâm ảo tạo thành đồ thị 14 node.
Trong pipeline hoàn toàn tự động này, độ chính xác giảm xuống \textbf{73,82\%}
(macro F1 = 0,7239), cho thấy sự phụ thuộc chặt chẽ giữa chất lượng ước lượng
tư thế và hiệu năng phân loại downstream. Phân tích theo từng lớp cho thấy các
hành động có chuyển động đặc trưng mạnh như \textit{jumping\_jacks} và
\textit{pullup} duy trì hiệu năng tốt (F1 $\geq$ 0,93), trong khi các bài tập
thực hiện ở tư thế nằm như \textit{bench\_press} suy giảm đáng kể (F1 = 0,38).
\end{abstract}

\begin{IEEEkeywords}
nhận dạng hành động, khung xương, ST-GCN, YOLOv8, ước lượng tư thế, Penn Action
\end{IEEEkeywords}

% ─────────────────────────────────────────────────────────────────────────────
\section{Giới thiệu}
\label{sec:intro}

Nhận dạng hành động người (Human Action Recognition -- HAR) là một bài toán cơ bản
trong thị giác máy tính với nhiều ứng dụng thực tiễn như giám sát an ninh, phân
tích thể thao, hỗ trợ phục hồi chức năng và tương tác người-máy. Các phương pháp
truyền thống dựa trên đặc trưng RGB thô nhạy cảm với sự thay đổi ngoại cảnh như
ánh sáng, trang phục và nền. Các phương pháp dựa trên khung xương~\cite{yan2018stgcn}
khắc phục hạn chế này bằng cách làm việc trên biểu diễn khớp cơ thể gọn nhẹ,
bất biến với góc nhìn.

Hai thách thức cản trở việc ứng dụng thực tế của HAR dựa trên khung xương là:
(1)~thu thập chuỗi khung xương chính xác mà không cần cảm biến chuyên dụng đắt
tiền hay nhân công gán nhãn; và (2)~xây dựng bộ phân loại bền vững trên sự biến
đổi thời gian đa dạng của các hành động. Các tiến bộ gần đây trong ước lượng
keypoint -- đặc biệt là họ YOLO -- cung cấp giải pháp thực tế cho thách thức đầu
tiên nhờ ước lượng tư thế một giai đoạn, thời gian thực từ ảnh RGB đơn kênh.
ST-GCN giải quyết thách thức thứ hai bằng cách mô hình hóa đồng thời cấu trúc
không gian của cơ thể người và sự tiến hóa thời gian của quỹ đạo các khớp.

Bài báo có các đóng góp sau:
\begin{itemize}
  \item Thiết kế pipeline HAR đầu cuối kết hợp YOLOv8-pose và ST-GCN, chỉ yêu
        cầu camera RGB thông thường.
  \item Thiết lập ngưỡng trên bằng cách huấn luyện ST-GCN trên keypoint ground
        truth của Penn Action, đạt 98,28\% độ chính xác.
  \item Định lượng khoảng cách hiệu năng khi dùng YOLO tự động (73,82\%) và phân
        tích nguyên nhân suy giảm theo từng lớp hành động.
  \item Cung cấp phân tích chất lượng keypoint YOLO theo từng khớp so với ground
        truth để giải thích suy giảm hiệu năng.
\end{itemize}

\begin{figure*}[t]
\centering
\resizebox{\textwidth}{!}{%
\begin{tikzpicture}[
  node distance=1.4cm,
  iobox/.style={trapezium, trapezium left angle=80, trapezium right angle=100,
                draw, thick, text width=2.0cm, align=center,
                font=\normalsize\bfseries, minimum height=1.6cm, inner sep=5pt},
  procbox/.style={rectangle, rounded corners=8pt, draw, thick,
                  text width=2.2cm, align=center,
                  font=\normalsize\bfseries, minimum height=1.6cm, inner sep=5pt},
  ann/.style={font=\small, text=black!55, align=center, text width=2.4cm},
  arr/.style={-Stealth, line width=2pt, color=black!45},
  dim/.style={font=\small\ttfamily, fill=white, inner sep=2pt,
              text=black!65, midway, above=2pt}
]
\node[iobox, fill=blue!18, draw=blue!55]
  (vid) {Video\\đầu vào};
\node[procbox, fill=orange!28, draw=orange!65, right=of vid]
  (yolo) {YOLOv8\\-pose};
\node[procbox, fill=teal!22, draw=teal!60, right=of yolo]
  (remap) {Ánh xạ\\khớp\\$17{\to}14$};
\node[procbox, fill=violet!18, draw=violet!55, right=of remap]
  (align) {Căn chỉnh\\thời gian\\$T{\to}64$};
\node[procbox, fill=red!18, draw=red!55, right=of align]
  (stgcn) {ST-GCN\\6 blocks\\GAP+FCN};
\node[iobox, fill=green!25, draw=green!55, right=of stgcn]
  (out) {Lớp\\hành động};

\draw[arr] (vid)   -- node[dim]{$(T,17,2)$}         (yolo);
\draw[arr] (yolo)  -- node[dim]{$(T,13,2)$}         (remap);
\draw[arr] (remap) -- node[dim]{$(T,14,2)$}         (align);
\draw[arr] (align) -- node[dim]{$(64,14,2)$}        (stgcn);
\draw[arr] (stgcn) -- node[dim]{$\mathrm{softmax}(8)$} (out);

\node[ann, below=0.3cm of vid]   {RGB frames};
\node[ann, below=0.3cm of yolo]  {17 COCO\\keypoints/frame};
\node[ann, below=0.3cm of remap] {Bỏ mắt/tai\\+ virtual center};
\node[ann, below=0.3cm of align] {Uniform\\index sampling};
\node[ann, below=0.3cm of stgcn] {$(N,2,64,14,1){\to}(N,8)$};
\node[ann, below=0.3cm of out]   {argmax\\(8 lớp)};
\end{tikzpicture}%
}
\caption{Pipeline đầu cuối nhận dạng hành động người đề xuất. Mỗi bước biến đổi
dữ liệu từ frame RGB thô đến nhãn hành động cuối cùng; nhãn trên mũi tên thể hiện
chiều tensor tại mỗi giai đoạn.}
\label{fig:pipeline}
\end{figure*}

% ─────────────────────────────────────────────────────────────────────────────
\section{Các Nghiên Cứu Liên Quan}
\label{sec:related}

\subsection{Nhận Dạng Hành Động Dựa Trên Khung Xương}
Yan \textit{và cộng sự}~\cite{yan2018stgcn} đề xuất ST-GCN, biểu diễn cơ thể
người dưới dạng đồ thị không-thời gian trong đó các node là khớp xương và các
cạnh mã hóa kết nối giải phẫu và chuyển tiếp thời gian. Mô hình phân vùng hàng
xóm thành ba tập con -- self-loop, hướng tâm (centripetal) và ly tâm (centrifugal)
-- cho phép học đặc trưng không gian thích nghi. ST-GCN trở thành baseline chuẩn
cho HAR dựa trên đồ thị và truyền cảm hứng cho nhiều công trình kế tiếp như
AGCN~\cite{shi2019agcn} và MS-G3D~\cite{liu2020msg3d}.

\subsection{Ước Lượng Tư Thế Phục Vụ HAR}
OpenPose~\cite{cao2019openpose} và HRNet~\cite{sun2019hrnet} cho chất lượng khung
xương cao nhưng yêu cầu pipeline đa giai đoạn tốn kém về tính toán. YOLOv8-pose
~\cite{ultralytics2023yolov8} tích hợp phát hiện đối tượng và hồi quy keypoint
trong một lần truyền mạng duy nhất, phù hợp cho ứng dụng thời gian thực. Một số
nghiên cứu gần đây đã sử dụng khung xương từ YOLO làm đầu vào cho bộ phân loại
GCN~\cite{liu2023yologcn}, cho thấy hiệu năng cạnh tranh với độ trễ thấp hơn
đáng kể.

\subsection{Penn Action Dataset}
Zhang \textit{và cộng sự}~\cite{zhang2013pennaction} giới thiệu Penn Action, tập
dữ liệu gồm 2.326 video clip thuộc 15 thể loại thể thao và bài tập với khung
xương 13 khớp được gán nhãn thủ công theo từng frame. Sự biến đổi cao trong cùng
một lớp (độ dài clip từ 18 đến 663 frame) và góc nhìn đa dạng làm cho đây trở
thành benchmark thử thách cho cả ước lượng tư thế lẫn phân loại hành động.

% ─────────────────────────────────────────────────────────────────────────────
\section{Pipeline Đề Xuất}
\label{sec:method}

\subsection{Tổng Quan Hệ Thống}

Hình~\ref{fig:pipeline} minh họa pipeline đầu cuối được đề xuất. Mỗi video đầu
vào được giải mã thành các frame riêng lẻ; YOLOv8-pose ước lượng khung xương 17
khớp COCO trên mỗi frame; bước ánh xạ chọn 13 khớp Penn Action và thêm một node
trung tâm ảo, tạo ra chuỗi khung xương 14 node; căn chỉnh thời gian lấy mẫu lại
chuỗi về đúng 64 frame; cuối cùng ST-GCN phân loại đồ thị không-thời gian thu
được thành một trong tám hành động bài tập.

\subsection{Trích Xuất Đặc Trưng Với YOLOv8-pose}
\label{sec:yolo}

Chúng tôi sử dụng biến thể \textit{yolov8n-pose} (nano, $\sim$6,5 MB), trả về
17 keypoint định dạng COCO trên mỗi người được phát hiện. Bộ keypoint COCO bao
gồm các điểm mặt (mắt, tai) không có đối ứng trong bộ gán nhãn 13 khớp của Penn
Action. Do đó, chúng tôi chỉ giữ lại 13 khớp có ý nghĩa cấu trúc bằng cách ánh
xạ chỉ số:
\[
  \text{COCO}_{[0,\,5{:}17]} \;\longrightarrow\; \text{Penn Action}_{[0{:}13]},
\]
bỏ qua chỉ số COCO 1--4 (mắt trái/phải và tai trái/phải). Khi không phát hiện
người trong một frame, toạ độ keypoint tương ứng được đặt về 0.

\subsection{Ánh Xạ Khớp Và Node Trung Tâm Ảo}
\label{sec:remap}

ST-GCN yêu cầu cấu trúc khung xương cố định. Theo cấu hình huấn luyện, chúng tôi
bổ sung \textbf{node trung tâm ảo} (khớp 13) được định nghĩa là tâm của bốn khớp
thân người:
\[
  \mathbf{p}_{13} = \frac{1}{4}\bigl(\mathbf{p}_{\text{l\_sho}} +
  \mathbf{p}_{\text{r\_sho}} + \mathbf{p}_{\text{l\_hip}} +
  \mathbf{p}_{\text{r\_hip}}\bigr).
\]
Node trung tâm này đóng vai trò gốc trong phân vùng đồ thị không gian
(xem Phần~\ref{sec:graph}), cung cấp điểm tham chiếu ổn định cho tổng hợp
hướng tâm và ly tâm.

\subsection{Căn Chỉnh Thời Gian}
\label{sec:temporal}

Độ dài clip Penn Action rất đa dạng (18--663 frame, trung bình 83,3). ST-GCN yêu
cầu chiều thời gian cố định. Chúng tôi sử dụng \textbf{lấy mẫu lại theo chỉ số
đều}: với clip có $T$ frame, chọn $T_{\text{target}} = 64$ chỉ số cách đều nhau
trên $[0, T-1]$ và thu thập keypoint tương ứng. Cách này giữ nguyên cấu trúc
chuyển động tổng thể, đồng thời bỏ qua frame dư trong clip dài và tránh lặp lại
frame biên.

\subsection{Kiến Trúc ST-GCN}
\label{sec:stgcn}

\subsubsection{Cấu Hình Đồ Thị Không Gian}
\label{sec:graph}
Khung xương được mô hình hóa là đồ thị $\mathcal{G} = (\mathcal{V}, \mathcal{E})$
trong đó $\mathcal{V}$ gồm 14 node và $\mathcal{E}$ gồm 17 cạnh giải phẫu (ví dụ:
vai--khuỷu tay, hông--đầu gối) cùng 4 cạnh bổ sung nối node trung tâm với hai
vai và hai hông. Theo ST-GCN~\cite{yan2018stgcn}, các node được phân vùng thành
ba tập so với khoảng cách đến node trung tâm:
\begin{itemize}
  \item $\mathbf{A}^{(0)}$: self-loop và hàng xóm cùng khoảng cách,
  \item $\mathbf{A}^{(1)}$: cạnh hướng tâm (node con $\to$ node cha),
  \item $\mathbf{A}^{(2)}$: cạnh ly tâm (node cha $\to$ node con).
\end{itemize}
Mỗi ma trận kề được chuẩn hóa theo bậc. Hình~\ref{fig:graph} minh họa sự phân
công phân vùng.

\begin{figure}[htbp]
  \centering
  \includegraphics[width=0.75\columnwidth]{charts/stgcn_notebook/graph_skeleton_partitions.png}
  \caption{Đồ thị khung xương 14 node tô màu theo loại phân vùng. Xanh lá:
  self-loop; cam: hướng tâm; đỏ: ly tâm. Node vàng (trung tâm) kết nối với
  hai vai và hai hông.}
  \label{fig:graph}
\end{figure}

\subsubsection{Kiến Trúc Mạng}
Mô hình xếp chồng sáu block ST-GCN. Mỗi block thực hiện tích chập đồ thị không
gian (GCN) tiếp theo là tích chập thời gian (TCN) với kernel $9 \times 1$ và
padding $4 \times 0$, cho trường tiếp nhận thời gian 9 frame. Bảng~\ref{tab:arch}
tổng hợp cấu hình.

\begin{table}[!t]
\caption{Cấu Hình Các Block ST-GCN}
\label{tab:arch}
\centering
\begin{tabular}{ccccc}
\toprule
Block & Kênh vào & Kênh ra & Stride & Residual \\
\midrule
1 & 2   & 64  & 1 & Không \\
2 & 64  & 64  & 1 & Có    \\
3 & 64  & 128 & 2 & Có    \\
4 & 128 & 128 & 1 & Có    \\
5 & 128 & 256 & 2 & Có    \\
6 & 256 & 256 & 1 & Có    \\
\bottomrule
\end{tabular}
\end{table}

Sau sáu block, global average pooling rút gọn chiều thời gian về 1, và một lớp
tích chập $1 \times 1$ ánh xạ vector đặc trưng 256 chiều thành 8 điểm số lớp.
Batch normalization đầu vào được áp dụng trên chiều $C \times V = 28$. Tensor đầu
vào có dạng $(N, C{=}2, T{=}64, V{=}14, M{=}1)$ trong đó $M$ là số người.

% ─────────────────────────────────────────────────────────────────────────────
\section{Thực Nghiệm}
\label{sec:exp}

\subsection{Dữ Liệu Và Thiết Lập}
\label{sec:setup}

\textbf{Tập dữ liệu.} Chúng tôi sử dụng Penn Action Dataset~\cite{zhang2013pennaction},
lọc lấy 8 loại bài tập tạo thành nhóm có ý nghĩa đồng nhất về mặt ngữ nghĩa:
\textit{bench\_press}, \textit{clean\_and\_jerk}, \textit{jump\_rope},
\textit{jumping\_jacks}, \textit{pullup}, \textit{pushup}, \textit{situp},
\textit{squat}. Tập con kết quả gồm \textbf{1.163 video}. Bảng~\ref{tab:dataset}
chi tiết phân phối theo lớp. Chúng tôi chia 80/20 cho train/validation (stratified,
\texttt{random\_state=42}), thu được 930 video huấn luyện và 233 video kiểm định.

\begin{table}[!t]
\caption{Phân Phối Tập Con Bài Tập Penn Action}
\label{tab:dataset}
\centering
\begin{tabular}{lcc}
\toprule
Hành động & Tổng số & Val \\
\midrule
squat            & 231 & 46 \\
pushup           & 211 & 42 \\
pullup           & 199 & 40 \\
bench\_press     & 140 & 28 \\
jumping\_jacks   & 112 & 23 \\
situp            & 100 & 20 \\
clean\_and\_jerk &  88 & 18 \\
jump\_rope       &  82 & 16 \\
\midrule
\textbf{Tổng}    & \textbf{1.163} & \textbf{233} \\
\bottomrule
\end{tabular}
\end{table}

Hình~\ref{fig:datastats} cho thấy tập con có sự mất cân bằng nhẹ giữa các lớp
(squat gấp gần 3 lần jump\_rope) và phân phối độ dài clip lệch phải mạnh, với
phần lớn video dưới 150 frame nhưng một số clip dài tới 663 frame.

\begin{figure}[!t]
  \centering
  \includegraphics[width=\columnwidth]{charts/stgcn_notebook/data_stats_and_length_distribution.png}
  \caption{Phân phối tập con bài tập Penn Action. Trái: số lượng video theo lớp
  (mất cân bằng nhẹ). Phải: phân phối độ dài clip thô trước khi căn chỉnh thời
  gian (lệch phải, mean $\approx$ 83 frame).}
  \label{fig:datastats}
\end{figure}

\textbf{Thiết lập huấn luyện.} ST-GCN được huấn luyện trên Google Colab với GPU
NVIDIA T4 bằng PyTorch. Optimizer: Adam ($lr = 10^{-3}$, weight decay $= 10^{-4}$);
learning rate scheduler: StepLR (step 30, $\gamma = 0,1$); batch size: 32; số
epoch: 80; hàm mất mát: cross-entropy.

\textbf{Chỉ số đánh giá.} Chúng tôi báo cáo độ chính xác Top-1, macro F1-score,
precision/recall/F1 theo từng lớp và confusion matrix.

% ── IV.B ──────────────────────────────────────────────────────────────────────
\subsection{Thực Nghiệm 1: ST-GCN Với Keypoint Ground Truth}
\label{sec:exp_gt}

Trong thực nghiệm đầu tiên, ST-GCN được huấn luyện và đánh giá sử dụng tọa độ
13 khớp được gán nhãn từ file \texttt{.mat} của Penn Action (cộng thêm node trung
tâm ảo), cung cấp ngưỡng trên cho hiệu năng phân loại khi có ước lượng tư thế
hoàn hảo.

\textbf{Quá trình huấn luyện.} Hình~\ref{fig:training} thể hiện các đường cong
huấn luyện. Mô hình hội tụ nhanh: đến epoch 19, độ chính xác validation vượt
93,99\%. Việc giảm learning rate tại epoch 30 (nhân $0,1$) đẩy nhanh hội tụ, và
mô hình đạt hiệu năng gần hoàn hảo ($\approx$97--99\%) từ epoch 31 trở đi.
Checkpoint tốt nhất (epoch 45) đạt 99,14\% accuracy và macro F1 = 0,9890.

\begin{figure*}[t]
  \centering
  \includegraphics[width=\textwidth]{charts/stgcn_notebook/training_curves.png}
  \caption{Lịch sử huấn luyện qua 80 epoch: loss (trái), accuracy (giữa) và macro
  F1 (phải) cho tập train (xanh) và validation (cam). Đường thẳng đứt nét đánh
  dấu thời điểm giảm learning rate tại epoch 30.}
  \label{fig:training}
\end{figure*}

\textbf{Kết quả.} Tại epoch 80, mô hình đạt \textbf{98,28\%} độ chính xác và
\textbf{0,9814} macro F1 trên 233 video validation. Tất cả tám lớp đều đạt F1
$\geq$ 0,96. Kết quả chi tiết theo từng lớp được đặt cạnh kết quả pipeline tại
Phần~\ref{sec:comparison} để tiện so sánh trực tiếp.

% ── IV.C ──────────────────────────────────────────────────────────────────────
\subsection{Thực Nghiệm 2: Pipeline Đầu Cuối YOLO + ST-GCN}
\label{sec:exp_pipeline}

Trong thực nghiệm thứ hai, keypoint ground truth được thay bằng keypoint từ
YOLOv8n-pose. Trọng số ST-GCN \textit{giữ nguyên} từ Thực nghiệm~1. Hình~\ref{fig:demo} minh
họa đầu ra điển hình của pipeline: khung xương 14 khớp từ YOLO overlay lên frame
và biểu đồ softmax confidence của ST-GCN.

\begin{figure*}[t]
  \centering
  \includegraphics[width=\textwidth]{charts/yolo_stgcn_notebook/yolo_skeleton_demo.png}
  \caption{Đầu ra pipeline trên video \texttt{0341} (\textit{bench\_press}). Trái:
  khung xương YOLO 14 khớp overlay lên frame; node vàng là trung tâm ảo được tính
  từ vai và hông. Phải: xác suất softmax ST-GCN trên 8 lớp.}
  \label{fig:demo}
\end{figure*}

Pipeline đạt \textbf{73,82\%} độ chính xác và \textbf{0,7239} macro F1, giảm
24,46 điểm phần trăm so với Thực nghiệm~1. Mức suy giảm không đồng đều giữa
các lớp:
\begin{itemize}
  \item \textit{jumping\_jacks} (F1 = 0,94) và \textit{pullup} (F1 = 0,93) duy
        trì hiệu năng cao do chuyển động biên độ lớn, đối xứng, ít bị ảnh hưởng
        bởi sai số keypoint.
  \item \textit{bench\_press} suy giảm nặng nhất (F1 = 0,38, recall = 25\%): tư
        thế nằm ngang khiến YOLO ước lượng sai và khung xương dễ nhầm với
        \textit{pushup} hoặc \textit{situp}.
  \item \textit{situp} bị nhiều lớp khác ``tràn vào'' (precision = 0,36), nhất
        quán với mẫu nhầm lẫn từ \textit{bench\_press}.
\end{itemize}

% ── IV.E ──────────────────────────────────────────────────────────────────────
\subsection{Tổng Hợp Và So Sánh}
\label{sec:comparison}

Bảng~\ref{tab:comparison_results} đặt kết quả hai thực nghiệm cạnh nhau.
Hình~\ref{fig:cm_compare} so sánh trực tiếp hai confusion matrix, làm rõ các
mẫu nhầm lẫn phát sinh khi chuyển từ GT keypoints sang YOLO keypoints.
Bảng~\ref{tab:summary} tóm tắt các chỉ số tổng thể.

\begin{table*}[!t]
\caption{So Sánh Kết Quả Phân Loại: ST-GCN + GT Keypoint (trái) Và YOLO + ST-GCN Pipeline (phải)}
\label{tab:comparison_results}
\centering
\begin{minipage}{0.47\textwidth}
\centering\small
\begin{tabular}{lcccc}
\toprule
Lớp & Prec. & Rec. & F1 & Sup. \\
\midrule
bench\_press     & 1,00 & 0,96 & 0,98 & 28 \\
clean\_and\_jerk & 0,95 & 1,00 & 0,97 & 18 \\
jump\_rope       & 1,00 & 0,94 & 0,97 & 16 \\
jumping\_jacks   & 1,00 & 1,00 & 1,00 & 23 \\
pullup           & 1,00 & 1,00 & 1,00 & 40 \\
pushup           & 0,95 & 0,98 & 0,96 & 42 \\
situp            & 1,00 & 0,95 & 0,97 & 20 \\
squat            & 0,98 & 1,00 & 0,99 & 46 \\
\midrule
\textbf{Accuracy} & \multicolumn{4}{c}{\textbf{98,28\%}} \\
\textbf{Macro F1} & \multicolumn{4}{c}{\textbf{0,9814}} \\
\bottomrule
\end{tabular}
\end{minipage}
\hfill
\begin{minipage}{0.47\textwidth}
\centering\small
\begin{tabular}{lcccc}
\toprule
Lớp & Prec. & Rec. & F1 & Sup. \\
\midrule
bench\_press     & 0,78 & 0,25 & 0,38 & 28 \\
clean\_and\_jerk & 0,75 & 0,83 & 0,79 & 18 \\
jump\_rope       & 0,70 & 0,88 & 0,78 & 16 \\
jumping\_jacks   & 0,92 & 0,96 & 0,94 & 23 \\
pullup           & 1,00 & 0,88 & 0,93 & 40 \\
pushup           & 0,76 & 0,60 & 0,67 & 42 \\
situp            & 0,36 & 0,80 & 0,49 & 20 \\
squat            & 0,81 & 0,83 & 0,82 & 46 \\
\midrule
\textbf{Accuracy} & \multicolumn{4}{c}{\textbf{73,82\%}} \\
\textbf{Macro F1} & \multicolumn{4}{c}{\textbf{0,7239}} \\
\bottomrule
\end{tabular}
\end{minipage}
\end{table*}

\begin{figure*}[!t]
\centering
\begin{minipage}{0.48\textwidth}
  \centering
  \includegraphics[width=\textwidth]{charts/stgcn_notebook/confusion_matrix_gt.png}
  \par\smallskip\small\textit{(a) ST-GCN + Ground Truth --- Accuracy = 98,28\%}
\end{minipage}
\hfill
\begin{minipage}{0.48\textwidth}
  \centering
  \includegraphics[width=\textwidth]{charts/yolo_stgcn_notebook/pipeline_confusion_matrix.png}
  \par\smallskip\small\textit{(b) YOLO + ST-GCN Pipeline --- Accuracy = 73,82\%}
\end{minipage}
\caption{So sánh confusion matrix giữa hai thực nghiệm trên cùng tập val 233
video. Với GT keypoints (a), hầu hết phần tử ngoài đường chéo bằng không. Với
YOLO keypoints (b), nhầm lẫn tập trung ở \textit{bench\_press} (bị nhầm sang
\textit{situp}/\textit{pushup}) và \textit{situp} (bị nhiễm từ lớp khác).}
\label{fig:cm_compare}
\end{figure*}

\begin{table}[!t]
\caption{Tổng Hợp Hiệu Năng}
\label{tab:summary}
\centering
\begin{tabular}{lcc}
\toprule
Cấu hình & Accuracy & Macro F1 \\
\midrule
ST-GCN + GT keypoint     & \textbf{98,28\%} & \textbf{0,9814} \\
YOLO + ST-GCN (pipeline) & 73,82\%          & 0,7239 \\
\midrule
Khoảng cách & $-$24,46\% & $-$0,2575 \\
\bottomrule
\end{tabular}
\end{table}

ST-GCN chứng tỏ là bộ phân loại rất hiệu quả khi được cung cấp khung xương chính
xác. Nút thắt cổ chai của hệ thống đầu cuối nằm ở chất lượng keypoint ước lượng
tự động, không phải ở năng lực của bộ phân loại.

% ─────────────────────────────────────────────────────────────────────────────
\section{Ablation Study: Chất Lượng Keypoint YOLO}
\label{sec:ablation}

Phần thực nghiệm cho thấy khoảng cách 24,46\% accuracy giữa hai cấu hình. Để
làm rõ nguồn gốc của sự suy giảm này, chúng tôi tiến hành phân tích định lượng
và định tính chất lượng keypoint YOLO so với ground truth Penn Action, đo trên
1.158 video bài tập (5 video bỏ qua do không phát hiện người), sai số chuẩn hóa
theo chiều cao người.

\begin{figure*}[t]
\centering
\begin{minipage}{0.48\textwidth}
  \centering
  \includegraphics[width=\textwidth]{charts/yolo_stgcn_notebook/keypoint_comparison_pushup.jpg}
  \par\smallskip\small\textit{(a) Video \#1348 (pushup): GT (đỏ) và YOLO (xanh lam)}
\end{minipage}
\hfill
\begin{minipage}{0.48\textwidth}
  \centering
  \includegraphics[width=\textwidth]{charts/yolo_stgcn_notebook/yolo_keypoint_quality.png}
  \par\smallskip\small\textit{(b) Sai số chuẩn hóa trung bình theo từng khớp (1.158 video)}
\end{minipage}
\caption{Phân tích chất lượng keypoint YOLO so với ground truth Penn Action.
Hình (a) minh họa sự lệch điển hình: YOLO đặt nhiều khớp chi dưới (đầu gối,
mắt cá chân) và khớp hông sai vị trí đáng kể so với annotation. Hình (b)
xác nhận định lượng: mắt cá chân có sai số cao nhất; đầu và vai tương đối ổn định.}
\label{fig:kp_ablation}
\end{figure*}

Sai số keypoint cao xuất phát từ ba nguyên nhân chính. (1)~\textbf{Domain
shift:} YOLOv8n-pose được huấn luyện trên COCO với đa dạng tỉ lệ người, trong
khi Penn Action quay người ở cự ly gần -- đầu hồi quy keypoint phải ngoại suy
ngoài phân phối huấn luyện. (2)~\textbf{Sai lệch định nghĩa khớp:} COCO và
Penn Action dùng điểm bám giải phẫu khác nhau (ví dụ: trung tâm hông COCO
so với mào chậu Penn Action), tạo ra lệch hệ thống không thể loại bằng ánh
xạ chỉ số. (3)~\textbf{Che khuất tư thế bài tập:} các tư thế như nằm (bench
press) hay ngồi xổm (squat) che khuất nghiêm trọng khớp chi dưới -- đúng là
nhóm có sai số phân loại lớn nhất.

% ─────────────────────────────────────────────────────────────────────────────
\section{Thảo Luận}
\label{sec:discussion}

\textbf{Độ bền vững và độ nhạy cảm.} Các hành động có chuyển động khớp biên độ
lớn, đối xứng, tần số cao (\textit{jumping\_jacks}, \textit{pullup}) phần lớn
không nhạy với nhiễu keypoint mức vừa phải. Ngược lại, các hành động thực hiện ở
góc nhìn không chuẩn hoặc bị che khuất một phần (\textit{bench\_press},
\textit{situp}) rất dễ bị ảnh hưởng. Điều này gợi ý các hướng cải thiện trong
tương lai: ước lượng tư thế bất biến góc nhìn, hoặc tăng cường dữ liệu với
keypoint nhiễu trong quá trình huấn luyện ST-GCN.

\textbf{Xử lý keypoint bị thiếu.} Khi YOLO không phát hiện người trong một frame,
chúng tôi đặt tất cả keypoint về 0. Điều này truyền một khung xương ``xẹp'' giả
tạo qua mô hình thời gian và có thể tạo ra các mẫu thời gian sai. Một giải pháp
hợp lý hơn -- như lấy frame hợp lệ gần nhất trước đó (forward-filling) hoặc nội
suy tuyến tính -- sẽ cải thiện độ bền vững.

\textbf{Căn chỉnh thời gian.} Lấy mẫu lại đều bỏ qua cấu trúc thời gian trong
clip rất dài (lên đến 663 frame, nén xuống 64). Chiến lược cửa sổ trượt thích
nghi (adaptive sliding window) hoặc pooling thời gian phân cấp có thể bảo toàn
tốt hơn các pha chuyển động có ý nghĩa.

\textbf{Quy mô mô hình YOLO.} Chúng tôi dùng \textit{yolov8n-pose} (nano, 6,5
MB) ưu tiên tốc độ. Các biến thể lớn hơn (\textit{yolov8x-pose}) có thể thu hẹp
khoảng cách chất lượng keypoint với chi phí độ trễ suy luận cao hơn.

% ─────────────────────────────────────────────────────────────────────────────
\section{Kết Luận}
\label{sec:conclusion}

Chúng tôi đã trình bày một pipeline nhận dạng hành động người đầu cuối kết hợp
YOLOv8-pose và ST-GCN, được đánh giá trên tập con 8 lớp bài tập của Penn Action.
ST-GCN với keypoint ground truth đạt 98,28\% độ chính xác, xác nhận năng lực biểu
diễn mạnh mẽ của mô hình cho phân loại hành động dựa trên khung xương. Pipeline
tự động hoàn toàn với keypoint từ YOLO đạt 73,82\%, khoảng cách $\approx$24,5
điểm phần trăm chủ yếu do sai số ước lượng tư thế -- đặc biệt ở các khớp mắt cá
chân và tư thế bài tập dễ bị che khuất.

Các kết quả này cho thấy \textit{chất lượng ước lượng tư thế là nút thắt hiệu
năng chính} trong pipeline HAR dựa trên khung xương. Các hướng nghiên cứu tiếp
theo bao gồm: (i)~tinh chỉnh (fine-tune) YOLOv8-pose trên dữ liệu miền bài tập
để giảm domain shift; (ii)~huấn luyện ST-GCN với tăng cường nhiễu keypoint để
cải thiện độ bền vững; và (iii)~khám phá tích chập đồ thị có trọng số tin cậy
(confidence-weighted) để giảm trọng số các khớp không chắc chắn trong suy luận.

% ─────────────────────────────────────────────────────────────────────────────
\begin{thebibliography}{00}

\bibitem{yan2018stgcn}
S. Yan, Y. Xiong, and D. Lin,
``Spatial Temporal Graph Convolutional Networks for Skeleton-Based Action Recognition,''
in \textit{Proc. AAAI Conf. Artif. Intell.}, 2018.

\bibitem{shi2019agcn}
L. Shi, Y. Zhang, J. Cheng, and H. Lu,
``Two-Stream Adaptive Graph Convolutional Networks for Skeleton-Based Action Recognition,''
in \textit{Proc. IEEE CVPR}, 2019.

\bibitem{liu2020msg3d}
Z. Liu, H. Zhang, Z. Chen, Z. Wang, and W. Ouyang,
``Disentangling and Unifying Graph Convolutions for Skeleton-Based Action Recognition,''
in \textit{Proc. IEEE CVPR}, 2020.

\bibitem{cao2019openpose}
Z. Cao, G. Hidalgo, T. Simon, S.-E. Wei, and Y. Sheikh,
``OpenPose: Realtime Multi-Person 2D Pose Estimation using Part Affinity Fields,''
\textit{IEEE Trans. Pattern Anal. Mach. Intell.}, vol. 43, no. 1, pp. 172--186, 2021.

\bibitem{sun2019hrnet}
K. Sun, B. Xiao, D. Liu, and J. Wang,
``Deep High-Resolution Representation Learning for Visual Recognition,''
\textit{IEEE Trans. Pattern Anal. Mach. Intell.}, vol. 43, no. 10, pp. 3349--3364, 2021.

\bibitem{ultralytics2023yolov8}
G. Jocher, A. Chaurasia, and J. Qiu,
``Ultralytics YOLOv8,'' 2023. [Online].
Available: \url{https://github.com/ultralytics/ultralytics}

\bibitem{liu2023yologcn}
Y. Liu, Z. Zhang, and C. Li,
``Real-Time Action Recognition Using YOLO-Based Pose Estimation and Graph Convolutional Networks,''
\textit{Sensors}, vol. 23, no. 8, p. 3912, 2023.

\bibitem{zhang2013pennaction}
W. Zhang, M. Zhu, and K. G. Derpanis,
``From Actemes to Action: A Strongly-supervised Representation for Detailed Action Understanding,''
in \textit{Proc. IEEE ICCV}, 2013.

\end{thebibliography}

\end{document}
